\documentclass{CarlosTexRpt}

\title{小型工程报告示例}
\author{Carlos QU}
\date{\today}

\begin{document}

\maketitle

\begin{abstract}
CarlosTexRpt 面向实验记录、设计说明、阶段总结和技术分析等小型中文工程报告。模板强调紧凑、清晰和低维护成本。
\end{abstract}

\tableofcontents

\chapter{设计目标}
这是章节标题后的第一段。该段应保持正常的首行缩进，用于验证标题后的正文排版行为已经修复。

这是同一章节中的第二段，用于对照普通段落的首行缩进。

\section{基本能力}
这是节标题后的第一段，同样应正常首行缩进。类文件内已显式设置标题后的缩进行为，不需要手工添加空命令或特殊占位。

\subsection{公式与单位}
系统采样频率设为 \SI{200}{\kilo\hertz}，采样周期为 \SI{5}{\micro\second}。示例公式为
\begin{equation}
  u = Ri + L\frac{\mathrm{d}i}{\mathrm{d}t}.
\end{equation}

\section{表格与代码}
表格采用三线表风格，如表~\ref{tab:params} 所示。

\begin{table}[H]
  \centering
  \caption{示例参数}
  \label{tab:params}
  \begin{tabular}{lll}
    \toprule
    参数 & 符号 & 数值 \\
    \midrule
    采样频率 & $f_s$ & \SI{200}{\kilo\hertz} \\
    母线电压 & $V_{dc}$ & \SI{24}{\volt} \\
    \bottomrule
  \end{tabular}
\end{table}

\begin{lstlisting}[language=json,caption={示例配置}]
{
  "sample_rate_hz": 200000,
  "dc_bus_v": 24
}
\end{lstlisting}

\end{document}
